\section{Conclusion}
\label{sec:conclusion}

We began by asking whether static analysis metrics could predict LLM-generated code correctness---a question motivated by the expensive alternative of running test suites on every candidate. Our work demonstrates that the answer is a strong yes: pylint score alone achieves $r=0.87$ correlation with pass@1 on HumanEval/MBPP, providing a cheap, effective quality signal.

\subsection{Summary}

This work makes three contributions:

\begin{enumerate}
\item \textbf{Empirical quantification:} We present the first correlation study establishing $r=0.87$ between pylint scores and functional correctness, far exceeding the $r\geq0.35$ threshold for moderate predictive power.

\item \textbf{Ensemble analysis:} We show that weighted combination of SA metrics actually degrades performance ($r=0.86 < 0.87$). This ``less is more'' finding simplifies deployment.

\item \textbf{Cross-model generalization:} All four tested LLMs show significant SA-correctness correlation (all $p<0.001$, all $r>0.35$), though with higher variance than anticipated.
\end{enumerate}

\subsection{Future Directions}

\textbf{From untested alternatives:} Non-linear ensemble methods (gradient boosting, neural combination) could discover synergies that simple weighting misses.

\textbf{From unverified assumptions:} Real API outputs from GPT-4, Claude, and other models would validate whether the GPT-4 outlier reflects true model differences or synthetic data artifacts.

\textbf{From scope extensions:} Repository-level benchmarks (SWE-bench) would test whether SA-correctness correlation extends to file-level and system-level code.

These findings open a practical avenue for LLM code deployment: cheap quality filtering that complements expensive test execution. We hope this work encourages further investigation of SA metrics as predictive signals, not just corrective feedback.
